% GENERATED FILE. Edit canonical lessons, then run npm run generate:books.
% canonical-source-hash: fnv1a64:71e6845f209b81d6
% canonical-lessons: TE-C322-duhkham, TE-S178-sign-visarga, TE-C322-bhayanga-undi, TE-C322-vellalani-undi, TE-R322-feelings-and-hopes

\chapter{Feelings and Hopes}
\label{ch:te-feelings-and-hopes}
\hlchaptermodality{322}

\begin{chapteropening}
\textbf{By the end of this chapter:} \emph{I can say how I feel and how I felt with nāku … -gā undi and anipiñcindi, and what I would like or dream of with -ālani.}

Complete TE-R322-feelings-and-hopes: say how you feel and felt, what you would like and dream of, and read the visarga \te{ః}.
\end{chapteropening}

\section[\texorpdfstring{\te{దుఃఖం} (duḥkhaṁ)}{దుఃఖం (duḥkhaṁ)}]{\te{దుఃఖం} (duḥkhaṁ) — grief}

\label{lesson:TE-C322-duhkham}



\begin{quote}
Before the new one: say \textquotedblleft{}sorrow, distress\textquotedblright{} (\emph{bādha}), \textquotedblleft{}comfort, ease\textquotedblright{} (\emph{sukhaṁ}) and \textquotedblleft{}a story\textquotedblright{} (\emph{katha}). Then \textquotedblleft{}my bag is lost\textquotedblright{}. (\emph{nā sañci pōyindi}.)
\end{quote}

\subsection*{You'll want to know: \te{దుఃఖం}}
\textbf{\te{దుఃఖం}} — \emph{duḥkhaṁ} — \textquotedblleft{}grief\textquotedblright{}.

\begin{quote}
\textbf{\te{ఆ} \te{కథ} \te{చదివి} \te{నాకు} \te{దుఃఖం} \te{వచ్చింది}.}
\end{quote}

\begin{quote}
\emph{ā katha cadivi nāku duḥkhaṁ vaccindi.} — \textquotedblleft{}Reading that story, I was moved to grief.\textquotedblright{}
\end{quote}

\textbf{\te{బాధ}} is the everyday hurt and sadness of a bad day. \textbf{\te{దుఃఖం}} is heavier: grief, the sorrow that brings tears. Like a fever, it \textquotedblleft{}comes\textquotedblright{} to you: \textbf{\te{నాకు} \te{దుఃఖం} \te{వచ్చింది}}.

The mark after \textbf{\te{దు}} is a breath, \emph{ḥ}. Many speakers let it run into the next sound, close to \emph{dukkhaṁ}. The next lesson is about that mark.

\begin{cousinweb}
\textbf{\te{దుఃఖం}} (\emph{duḥkhaṁ}) is borrowed from Sanskrit \emph{duḥkha}, \textquotedblleft{}difficulty, pain\textquotedblright{} (CDIAL 6375), with Telugu's noun ending \textbf{-\te{ం}}. CDIAL notes it was formed after \emph{sukha}, the word you know as \textbf{\te{సుఖం}}: the two are a pair, ease and its opposite.
\end{cousinweb}

\subsection*{Your turn}
Say these aloud:

\begin{itemize}
\raggedright
  \item \emph{duḥkhaṁ}, with a short breath after \emph{du}
  \item \textquotedblleft{}Reading that story, I was moved to grief\textquotedblright{} — \emph{ā katha cadivi nāku duḥkhaṁ vaccindi}
  \item the everyday word for sadness, then the heavier one
  \item the pair: \emph{sukhaṁ}, \emph{duḥkhaṁ}
\end{itemize}

\begin{tcolorbox}[breakable,colback=teal!4,colframe=teal!35!black,title={Before you move on}]
What does \textbf{\te{దుఃఖం}} mean? (\textbf{Grief.}) How is it different from \textbf{\te{బాధ}}? (\textbf{It is heavier.}) Which word is its opposite? (\textbf{\te{సుఖం}}.)
\end{tcolorbox}
\section[\texorpdfstring{\te{ః} (ḥ)}{ః (ḥ)}]{\te{ః} (ḥ) — the visarga inside \te{దుఃఖం}}

\label{lesson:TE-S178-sign-visarga}



\begin{quote}
Before the new one: write \textbf{\te{ం}} and \textbf{\te{ఖ}} once each, and name them. (The sunna, \emph{ṁ}, and \emph{kha}.)

Now say \textbf{\te{దుఃఖం}} (\emph{duḥkhaṁ}), \textquotedblleft{}grief\textquotedblright{}. Both of those are in it. This lesson takes out the one mark you have not written and puts it in your hand.
\end{quote}

\subsection*{Script you'll notice: \te{ః}}
\textbf{\te{ః}} — the visarga, \emph{ḥ}: a short breath after a vowel. It is two small rings, one above the other, written after its syllable: \textbf{\te{దు}} + \textbf{\te{ః}} $\to$ \textbf{\te{దుః}}, \emph{duḥ}.

Do not confuse it with \textbf{\te{ం}}, one ring, which is \emph{ṁ}. The visarga lives in words taken from Sanskrit, and most of them are learned words; \textbf{\te{దుఃఖం}} is the one you will meet most.

\subsection*{Writing: \te{ః} — copy the printed shape}
Write \textbf{\te{ః}} yourself, slowly, and larger than it is printed: two small rings the same size, one above the other, both on the right of the syllable they follow.

\begin{quote}
This book does not yet say where the pen starts on this shape, or which way it travels. That is taught with real variation from school to school, and it waits here until it can be written down from a source. Copying the printed shape needs no such source, and it puts the shape in your hand in the meantime.
\end{quote}

\subsection*{Your turn}
\begin{itemize}
\raggedright
  \item \emph{Find:} \te{ః} inside \textbf{\te{దుఃఖం}}
  \item \emph{Trace it:} \te{ః} three times, saying \emph{ḥ} as you finish each one
  \item \emph{Write it:} \te{ం}, then \te{ః}, from memory, naming each as you finish it
  \item \emph{Write it:} \textbf{\te{దుఃఖం}}
\end{itemize}

\begin{tcolorbox}[breakable,colback=teal!4,colframe=teal!35!black,title={Before you move on}]
Which mark is this — \te{ః}? (\textbf{The visarga}, \emph{ḥ}, a breath.) Which word did it come from? (\textbf{\te{దుఃఖం}}, \emph{duḥkhaṁ}, \textquotedblleft{}grief\textquotedblright{}.) How is it different from \textbf{\te{ం}}? (\textbf{Two rings, not one.})
\end{tcolorbox}
\section[\texorpdfstring{\te{నాకు} \te{భయంగా} \te{ఉంది} (nāku bhayaṅgā undi)}{నాకు భయంగా ఉంది (nāku bhayaṅgā undi)}]{\te{నాకు} \te{భయంగా} \te{ఉంది} (nāku bhayaṅgā undi) — how I feel}

\label{lesson:TE-C322-bhayanga-undi}



\begin{quote}
Say \textquotedblleft{}fear\textquotedblright{} (\emph{bhayaṁ}), \textquotedblleft{}anger\textquotedblright{} (\emph{kōpaṁ}) and \textquotedblleft{}happily\textquotedblright{} (\emph{santōṣaṅgā}). Then say \textquotedblleft{}grief\textquotedblright{}. (\emph{duḥkhaṁ}.) Which mark sits after its first syllable? (\textbf{\te{ః}}, the visarga, \emph{ḥ}.)
\end{quote}

\subsection*{You'll want to know: \te{నాకు} … -\te{గా} \te{ఉంది}}
\begin{quote}
\textbf{\te{నాకు} \te{చాలా} \te{సంతోషంగా} \te{ఉంది}.}
\end{quote}

\begin{quote}
\emph{nāku cālā santōṣaṅgā undi.} — \textquotedblleft{}I am very happy.\textquotedblright{}
\end{quote}

\begin{quote}
\textbf{\te{నాకు} \te{భయంగా} \te{ఉంది}.}
\end{quote}

\begin{quote}
\emph{nāku bhayaṅgā undi.} — \textquotedblleft{}I am afraid.\textquotedblright{}
\end{quote}

The one who feels is in the dative, \textbf{\te{నాకు}}, as in \textbf{\te{నాకు} \te{తెలుగు} \te{వచ్చు}}. The feeling takes \textbf{-\te{గా}}, and \textbf{\te{ఉంది}}, \textquotedblleft{}it is\textquotedblright{}, closes the sentence: \textquotedblleft{}to me, fearfully, it is\textquotedblright{}. So too \textbf{\te{నాకు} \te{బాధగా} \te{ఉంది}}, \textquotedblleft{}I feel sad\textquotedblright{}; \textbf{\te{నాకు} \te{చలిగా} \te{ఉంది}}, \textquotedblleft{}I feel cold\textquotedblright{}; \textbf{\te{నాకు} \te{ఆకలిగా} \te{ఉంది}}, \textquotedblleft{}I am hungry\textquotedblright{}.

Some feelings \textquotedblleft{}come\textquotedblright{} instead, as \textbf{\te{దుఃఖం}} does:

\begin{quote}
\textbf{\te{నాకు} \te{కోపం} \te{వచ్చింది}.}
\end{quote}

\begin{quote}
\emph{nāku kōpaṁ vaccindi.} — \textquotedblleft{}I got angry.\textquotedblright{}
\end{quote}

\subsection*{You'll want to know: … \te{అనిపించింది}, I felt}
\begin{quote}
\textbf{\te{రాత్రి} \te{ఒంటరిగా} \te{నడుస్తున్నప్పుడు} \te{నాకు} \te{భయంగా} \te{అనిపించింది}.}
\end{quote}

\begin{quote}
\emph{rātri oṇṭarigā naḍustunnappuḍu nāku bhayaṅgā anipiñcindi.} — \textquotedblleft{}Walking alone at night, I felt afraid.\textquotedblright{}
\end{quote}

For how a moment felt, put the \textbf{-\te{గా}} feeling before \textbf{\te{అనిపించింది}}, \textquotedblleft{}it seemed\textquotedblright{}: \textquotedblleft{}to me, fearfully, it seemed\textquotedblright{}. \textbf{\te{ఒంటరిగా}} (\emph{oṇṭarigā}, \textquotedblleft{}alone\textquotedblright{}). The scene is the \textbf{-\te{తున్నప్పుడు}} of a story; the feeling is the event in it.

\subsection*{Your turn}
Say these aloud:

\begin{itemize}
\raggedright
  \item \textquotedblleft{}I am very happy\textquotedblright{} — \emph{nāku cālā santōṣaṅgā undi}
  \item \textquotedblleft{}I am afraid\textquotedblright{} — \emph{nāku bhayaṅgā undi}
  \item \textquotedblleft{}I got angry\textquotedblright{} — \emph{nāku kōpaṁ vaccindi}
  \item \textquotedblleft{}Walking alone at night, I felt afraid\textquotedblright{} — \emph{rātri oṇṭarigā naḍustunnappuḍu nāku bhayaṅgā anipiñcindi}
  \item how you felt on your last trip, and when
\end{itemize}

\begin{tcolorbox}[breakable,colback=teal!4,colframe=teal!35!black,title={Before you move on}]
\textquotedblleft{}I feel sad\textquotedblright{}? (\textbf{\te{నాకు} \te{బాధగా} \te{ఉంది}}.) \textquotedblleft{}I am hungry\textquotedblright{}? (\textbf{\te{నాకు} \te{ఆకలిగా} \te{ఉంది}}.) \textquotedblleft{}I felt afraid\textquotedblright{}? (\textbf{\te{నాకు} \te{భయంగా} \te{అనిపించింది}}.) Which feeling \textquotedblleft{}came\textquotedblright{}? (\textbf{Anger:} \textbf{\te{నాకు} \te{కోపం} \te{వచ్చింది}}.)
\end{tcolorbox}
\section[\texorpdfstring{\te{వెళ్ళాలని} \te{ఉంది} (veḷḷālani undi)}{వెళ్ళాలని ఉంది (veḷḷālani undi)}]{\te{వెళ్ళాలని} \te{ఉంది} (veḷḷālani undi) — what I would like}

\label{lesson:TE-C322-vellalani-undi}



\begin{quote}
Say \textquotedblleft{}I must go\textquotedblright{} (\emph{nēnu veḷḷāli}) and \textquotedblleft{}I am happy\textquotedblright{} (\emph{nāku santōṣaṅgā undi}). Then \textquotedblleft{}a dream\textquotedblright{}, \textquotedblleft{}hope\textquotedblright{} and \textquotedblleft{}a wish\textquotedblright{}. (\emph{kala}, \emph{āśa}, \emph{kōrika}.)
\end{quote}

\subsection*{You'll want to know: \te{నాకు} … -\te{ఆలని} \te{ఉంది}}
\begin{quote}
\textbf{\te{నాకు} \te{తెలుగు} \te{బాగా} \te{నేర్చుకోవాలని} \te{ఉంది}.}
\end{quote}

\begin{quote}
\emph{nāku telugu bāgā nērcukōvālani undi.} — \textquotedblleft{}I would like to learn Telugu well.\textquotedblright{}
\end{quote}

\begin{quote}
\textbf{\te{నాకు} \te{మా} \te{ఊరికి} \te{వెళ్ళాలని} \te{ఉంది}.}
\end{quote}

\begin{quote}
\emph{nāku mā ūriki veḷḷālani undi.} — \textquotedblleft{}I would like to go to my home town.\textquotedblright{}
\end{quote}

Take the \textbf{-\te{ఆలి}} of \textquotedblleft{}must\textquotedblright{}, add \textbf{\te{అని}}, \textquotedblleft{}saying\textquotedblright{}: \textbf{\te{వెళ్ళాలి}} + \textbf{\te{అని}} $\to$ \textbf{\te{వెళ్ళాలని}}. Then \textbf{\te{నాకు} … \te{ఉంది}}: \textquotedblleft{}to me there is (the thought) that I must go\textquotedblright{}. The wish sits in the dative, like a feeling.

\subsection*{You'll want to know: my wish, my dream, my hope}
\begin{quote}
\textbf{\te{తెలుగులో} \te{ఒక} \te{పుస్తకం} \te{రాయాలని} \te{నా} \te{కోరిక}.}
\end{quote}

\begin{quote}
\emph{telugulō oka pustakaṁ rāyālani nā kōrika.} — \textquotedblleft{}My wish is to write a book in Telugu.\textquotedblright{}
\end{quote}

\begin{quote}
\textbf{\te{వైద్యుడు} \te{కావాలని} \te{నా} \te{కల}.}
\end{quote}

\begin{quote}
\emph{vaidyuḍu kāvālani nā kala.} — \textquotedblleft{}My dream is to become a doctor.\textquotedblright{}
\end{quote}

\begin{quote}
\textbf{\te{ఆ} \te{ఉద్యోగం} \te{వస్తుందని} \te{నా} \te{ఆశ}.}
\end{quote}

\begin{quote}
\emph{ā udyōgaṁ vastundani nā āśa.} — \textquotedblleft{}My hope is that I will get that job.\textquotedblright{}
\end{quote}

The wish comes first, closed by \textbf{-\te{ఆలని}} or \textbf{-\te{అని}}; \textbf{\te{నా} \te{కోరిక}}, \textbf{\te{నా} \te{కల}} or \textbf{\te{నా} \te{ఆశ}} comes last. In \textbf{\te{కావాలని}}, \textbf{\te{కా}-} has its plain sense, \textquotedblleft{}to become\textquotedblright{}. A job \textquotedblleft{}comes\textquotedblright{} to you: \textbf{\te{ఉద్యోగం} \te{వస్తుంది}}.

\subsection*{Your turn}
Say these aloud:

\begin{itemize}
\raggedright
  \item \textquotedblleft{}I would like to learn Telugu well\textquotedblright{} — \emph{nāku telugu bāgā nērcukōvālani undi}
  \item \textquotedblleft{}I would like to go to my home town\textquotedblright{} — \emph{nāku mā ūriki veḷḷālani undi}
  \item \textquotedblleft{}My wish is to write a book in Telugu\textquotedblright{} — \emph{telugulō oka pustakaṁ rāyālani nā kōrika}
  \item \textquotedblleft{}My dream is to become a doctor\textquotedblright{} — \emph{vaidyuḍu kāvālani nā kala}
  \item one wish and one dream of your own
\end{itemize}

\begin{tcolorbox}[breakable,colback=teal!4,colframe=teal!35!black,title={Before you move on}]
\textquotedblleft{}I would like to go\textquotedblright{}? (\textbf{\te{నాకు} \te{వెళ్ళాలని} \te{ఉంది}}.) Which two pieces make \textbf{-\te{ఆలని}}? (\textbf{-\te{ఆలి}} + \textbf{\te{అని}}.) \textquotedblleft{}My dream is …\textquotedblright{}? (\textbf{… -\te{ఆలని} \te{నా} \te{కల}}.)
\end{tcolorbox}
\section[(review)]{Review — feelings and hopes}

\label{lesson:TE-R322-feelings-and-hopes}



\begin{quote}
Say \textquotedblleft{}grief\textquotedblright{}, \textquotedblleft{}a remedy\textquotedblright{} and \textquotedblleft{}one's home place\textquotedblright{}. (\emph{duḥkhaṁ}, \emph{auṣadhaṁ}, \emph{ūru}.)

\begin{itemize}
\raggedright
  \item \emph{Write it:} \textbf{\te{ః}}, \textbf{\te{ఔ}} and \textbf{\te{ఊ}}, once each from memory
\end{itemize}
\end{quote}

\subsection*{Your turn}
\begin{itemize}
\raggedright
  \item \emph{Say it:} I am very happy. I would like to go to my home town.
  \item \emph{Say it:} When the bus went, I felt angry.
  \item \emph{Say it:} My bag was lost. I felt sad.
  \item \emph{Say it:} My dream is to write a book. My hope is that I will get that job.
  \item \emph{Write it:} \textbf{\te{దుఃఖం}}, \textbf{\te{ఔషధం}}, \textbf{\te{ఊరు}}, \textbf{\te{ఓటు}}, \textbf{\te{ఈలోగా}}
\end{itemize}

\begin{tcolorbox}[breakable,colback=teal!4,colframe=teal!35!black,title={Before you move on}]
\textquotedblleft{}I am afraid\textquotedblright{}? (\textbf{\te{నాకు} \te{భయంగా} \te{ఉంది}}.) \textquotedblleft{}I felt happy\textquotedblright{}? (\textbf{\te{నాకు} \te{సంతోషంగా} \te{అనిపించింది}}.) \textquotedblleft{}I would like to learn\textquotedblright{}? (\textbf{\te{నాకు} \te{నేర్చుకోవాలని} \te{ఉంది}}.) \textquotedblleft{}Grief\textquotedblright{}? (\textbf{\te{దుఃఖం}}, with the two rings of \textbf{\te{ః}}.)
\end{tcolorbox}
